\begin{tikzpicture}
\node[box,text width=2.6cm] (front) at (0,0)
 {接入前端\\认证／授权\\函数元数据};
\node[good,text width=3cm] (router) at (4.2,1.3)
 {路由管理者\\找到已有可用槽\\协调调用使用};
\node[box,text width=3cm] (place) at (4.2,-1.6)
 {放置服务\\无槽时选择 worker\\管理槽寿命与负载};
\node[good,text width=3.3cm] (worker) at (8.8,0)
 {worker\\MicroVM 执行槽\\运行时／函数代码\\先后调用复用};
\draw[ctl] (front)--(router);
\draw[ctl] (router)--node[right,note]{申请新槽} (place);
\draw[ctl] (place)--(worker);
\draw[ctl] (router)--(worker);
\draw[arr] (front.south)--(0,-3.2)--node[below,note]{前端取得位置后直接发送事件载荷} (8.8,-3.2)--(worker.south);
\node[note,text width=11.3cm] at (4.4,-4.3)
 {按 Firecracker 公开部署设计表示：虚线为控制，实线为事件数据。};
\end{tikzpicture}
